% ================================================================
\section{Banco de questões pós-AV1 e checklist}\label{sec:bancoav2}
% ================================================================
As listas do curso estão resolvidas nas próprias seções: PD em §\ref{sec:listapd}, AGM e gulosos em §\ref{sec:listaagm}. Aqui ficam as listas de outras disciplinas (pasta ``Outros Materiais'': APA/UFF, MAC5711/USP lista 6, MC458/Unicamp lista 7), que servem de treino extra. Os valores numéricos foram conferidos por programa.

\subsection{Listas extras: respostas e ideias}
\begin{center}\small
\begin{tabularx}{\linewidth}{@{}lX@{}}
\toprule
Exercício & Resposta / recorrência\\
\midrule
APA 1--5 & Backtracking (permutações, combinações, ciclo hamiltoniano, caixeiro viajante, damas). Fora do programa; 6 damas, primeira solução da busca linha a linha (colunas crescentes): $(2,4,6,1,3,5)$.\\
APA 6 & LCS de \textrm{dacdba} e \textrm{cadcda}: 4 (\textrm{acda} ou \textrm{dcda}) --- §\ref{sec:lcs}.\\
APA 7 & Mochila 0-1, $W=100$: 240, com os itens 1, 2 e 4 --- §\ref{sec:mochila}.\\
APA 8 & Kadane (PD em $O(n)$): $S(i)=\max\{a_i,\ S(i-1)+a_i\}$ é a melhor soma terminando em $i$; resposta $\max_iS(i)$. $5,15,-30,10,-5,40,10\to$ \textbf{55} ($10,-5,40,10$).\\
\midrule
l6 1 & Cadeia $\langle5,10,3,12,5,50,6\rangle$: \textbf{2010}, $((A_1A_2)((A_3A_4)(A_5A_6)))$ --- §\ref{sec:cadeia}.\\
l6 2, 3, 4 & Tabela $s$ e \textsc{Multiplica}; todas as LCS; maior subsequência crescente $O(n^2)$ e $O(n\log n)$ --- §\ref{sec:cadeia}, §\ref{sec:lcs}, §\ref{sec:listapd}.\\
l6 5 & Número de ABBs com 6 chaves: Catalan $C_6=\binom{12}{6}/7=$ \textbf{132}.\\
l6 6--10 & ABB ótima (Cormen §15.5, não dada em aula): $c(i,j)=\sum_{k=i}^{j}a_k+\min_{i\le r\le j}\{c(i,r-1)+c(r+1,j)\}$, $c(i,i-1)=0$, em $\Theta(n^3)$ (mesma forma da cadeia de matrizes). Ex.~6: frequências $1,1,2,3$: a raiz ótima é a chave 3 (custo 13; com a 4 na raiz, 14). Ex.~7: custo 185, raiz \emph{if}. Ex.~8: some também os $b$: $w(i,j)=\sum_{i}^{j}a_k+\sum_{i-1}^{j}b_k$, base $c(i,i-1)=b_{i-1}$. Ex.~9 (Cormen 15.5-2): custo \textbf{312} (3{,}12 em probabilidades), raiz $k_5$.\\
l6 11 & Subsequência crescente de soma máxima: $S[i]=a_i+\max\{S[j]:j<i,\,a_j<a_i\}$ (0 se vazio), $O(n^2)$; guarde $ant[i]$. Ex.: $1,101,2,3,100,4,5\to106$ ($1,2,3,100$).\\
l6 12 / MC458 8 & Impressão nítida: $c(j)=\min_{i\le j,\ \text{cabe}}\{c(i-1)+\text{folga}(i,j)^3\}$ (na l6 a última linha custa 0; na MC458 a penalidade é o quadrado e inclui a última), $O(n^2)$.\\
l6 13 & Festa da empresa (PD em árvore): $com(v)=s(v)+\sum_{f}sem(f)$; $sem(v)=\sum_f\max\{com(f),sem(f)\}$; resposta $\max$ na raiz. $\Theta(n)$ (pós-ordem).\\
l6 14 / MC458 10 & Cortes (tora/string): §\ref{sec:listapd}, ex.~6. Tora de 10 m com cortes $2,4,7$: \textbf{20}.\\
l6 15 & Torres gêmeas $=$ LCS das duas sequências de raios.\\
l6 20 & Bandeiras: $f(1)=f(2)=2$ e $f(n)=f(n-1)+f(n-2)$: a última faixa é sempre B ou V; se a penúltima não é azul, cada configuração válida de $n-1$ faixas se estende de um único jeito; se é azul, cada uma de $n-2$ faixas se estende de um único jeito (A $+$ a cor que falta). Valores $2,2,4,6,10,16,\dots$\\
\midrule
MC458 3 & Mochila exata (subset sum), pesos $2,7,5,2$, $K=0..14$; ``O'' $=$ há solução sem o item, ``I'' $=$ só usando o item: \newline
\texttt{p3=5: O - O - - I - O - O - - I - I}\newline
\texttt{p4=2: O - O - I O - O - O - I O - O}\\
MC458 4 & Memória $O(K)$: um vetor, percorrendo $K$ \emph{decrescendo} --- §\ref{sec:mochila}.\\
MC458 5 & Mochila ilimitada --- §\ref{sec:listapd}, ex.~5.\\
MC458 7 & Submatriz de soma máxima: fixe o par de linhas $(i_1,i_2)$ ($O(m^2)$ pares), some as colunas e rode Kadane ($O(n)$): $O(m^2n)$.\\
MC458 9 & Laje de mármore: $P(a,b)=\max\{p_{ab},\ \max_{a'<a}P(a',b)+P(a-a',b),\ \max_{b'<b}P(a,b')+P(a,b-b')\}$; $O(nm(n+m))\subseteq O(n^2m^2)$.\\
\bottomrule
\end{tabularx}
\end{center}

\subsection{Como estudar a parte nova}
O formato da AV2 \textbf{ainda não foi anunciado} nos materiais disponíveis. Pelo estilo da AV1 e das listas, as questões mais prováveis são:
\begin{prova}[Tipos de questão prováveis (baseado nas listas e no estilo da AV1)]
\begin{itemize}
  \item \textbf{PD}: dado um problema novo, escrever subproblema, recorrência, pseudocódigo e complexidade; preencher uma tabela pequena (LCS, cadeia de matrizes, mochila) e reconstruir a solução.
  \item \textbf{Simulação}: rodar BFS, DFS (com $PE/PS$, tipos de arestas, $back$/$old$), Dijkstra, Floyd-Warshall (matrizes $D^k$ e $P$), Kruskal e Prim num grafo dado.
  \item \textbf{Provas}: correção de Kruskal e Prim (propriedade do corte), do guloso em matroides, de Dijkstra; verificar se um sistema de conjuntos é matroide (provar (I3) ou dar contraexemplo).
  \item \textbf{Escolha de implementação}: Prim/Dijkstra com matriz ($O(n^2)$) ou com heap ($O(m\log n)$), conforme a densidade; Union-Find no Kruskal.
\end{itemize}
\end{prova}

\begin{chave}[Checklist da parte pós-AV1]
\begin{enumerate}
  \item \textbf{PD}: subestrutura ótima $+$ sobreposição; as três recorrências (LCS, cadeia, mochila) de cor, com as ordens de preenchimento; por que a mochila é pseudo-polinomial.
  \item \textbf{Floyd-Warshall}: significado de $D^k$, recorrência, $k$ externo, matriz $P$ (penúltimo vértice) e reconstrução.
  \item \textbf{Dijkstra}: pseudocódigo, tabela de execução, prova, contraexemplo negativo, as três complexidades.
  \item \textbf{Buscas}: pseudocódigos com $PE/PS$ e com níveis; tipos de arestas e suas caracterizações; aplicações (componentes, ciclos, bipartição, distâncias, articulações, CFC, ordenação topológica).
  \item \textbf{Gulosos e matroides}: axiomas; exemplos e não exemplos; teorema do guloso (com a prova por troca) e a recíproca.
  \item \textbf{AGM}: propriedade do corte; Kruskal (Union-Find, $O(m\log n)$) e Prim (as quatro implementações); AGM de peso máximo.
\end{enumerate}
\end{chave}
